\documentclass[../main]{subfiles}
\begin{document}
\ifSubfilesClassLoaded{\setcounter{page}{4}}{}
\section{L’opera necessaria e le condizioni per l’attività comunista}
\label{sec:06_opera_necessaria}

In una società comunista, l’esistenza dell’individuo non è più legata alla proprietà privata e allo scambio di merci, liberandolo così da tutte le necessità connesse alla salvaguardia delle condizioni di esistenza di tale proprietà. Questo offre alla società notevoli vantaggi, preservandone al contempo le risorse umane. Nei “Discorsi di Elberfeld”, Engels citò una serie di esempi che illustrano con chiarezza quanta forza lavoro umana potrebbe essere risparmiata in una società in cui gli interessi dei suoi membri non sono in conflitto e i frutti del lavoro non si trasformano in merci destinate a competere tra loro sul mercato.

L’impegno richiesto per il funzionamento del mercato, meccanismo attraverso il quale le risorse sociali vengono distribuite, cede il passo a una più efficiente allocazione diretta. Ciò consente di evitare attività prive di senso nell’ambito dell’organizzazione della produzione su scala sociale, risparmiando tempo prezioso che altrimenti verrebbe sprecato in incombenze inutili. Non avrà più ragione d’esistere la produzione di imballaggi superflui destinati alla mera conservazione dei prodotti, né i trasporti legati alle dinamiche di vendita piuttosto che alla consegna al consumatore finale; si eviterà la produzione di beni destinati a rimanere invenduti e a divenire eccedenze di magazzino, e, soprattutto, si porrà fine alla fabbricazione di prodotti dannosi per la salute, ma fonte di enormi profitti in un sistema economico basato sulla produzione di merci (droghe, e simili). Tutto ciò risulterà inconcepibile in una società in cui la produzione non è più orientata al profitto.

L’abolizione del carattere privato dell’appropriazione comporta, altresì, un notevole risparmio di risorse umane. Si tratta dello stesso principio che regola la produzione dei beni, ma analizzato alla luce dell’impiego di manodopera necessario per raggiungere la fase di consumo finale. Gli esempi tratti dalla vita quotidiana e dal settore dei trasporti, già menzionati, illustrano in modo inequivocabile la possibilità di liberare una considerevole quantità di lavoro umano, attualmente impiegata in attività inefficienti e frammentate.

Nell'era del capitalismo moderno, si esige che gli individui indirizzino le proprie energie verso il sostegno della capacità produttiva di plusvalore. Tale forma di attività lavorativa si diffonde sempre più. Parallelamente alla produzione di plusvalore, si è sviluppata ineluttabilmente anche la produzione di beni di consumo, dando origine a uno stile di vita incentrato sull'acquisizione di determinati prodotti. Questo stesso dato evidenzia le crescenti (e significativamente superiori rispetto all'epoca di Marx ed Engels) possibilità di emancipazione, in cui gli individui, operando in un contesto di proprietà collettiva, potranno dedicarsi alla produzione di ciò che è realmente necessario per il progresso della società.

Gran parte del tempo e delle risorse della società vengono impiegati per sostenere le restrizioni che la proprietà privata esercita sulla vita sociale e per attenuarne gli effetti più nefasti, inclusi i crimini. A tal fine, si ricorre a forze dell'ordine, esercito, personale amministrativo e a diverse tipologie di guardie addette alla protezione della proprietà privata. Liberando questa forza lavoro, attraverso il superamento della proprietà privata sui mezzi di produzione, e riindirizzandola verso attività di interesse diretto per la società, si otterrebbe una significativa riduzione del tempo di lavoro necessario e si creerebbe spazio per il tempo libero, aprendo la strada a una vita pienamente umana, ovvero al comunismo.

Questo cambiamento, di per sé, e sulla base delle tecnologie già disponibili, consentirebbe di ridurre sensibilmente l'orario di lavoro e di ampliare le opportunità per il tempo libero, per una vita veramente umana: il comunismo.

Un'altra significativa tendenza, che si sviluppa all'interno del sistema capitalistico e la cui piena realizzazione è ostacolata unicamente dalla sua natura di rapporto sociale predominante, prefigura l'avvento del comunismo.«L’importanza storica del capitale risiede nella capacità di generare un surplus di lavoro, un’eccedenza che trascende la mera necessità di sussistenza e il semplice valore d’uso. La sua funzione storica si esaurirà quando, da un lato, i bisogni umani si saranno evoluti al punto che questo surplus, il lavoro che va al di là dello strettamente indispensabile per la sopravvivenza, diverrà un’esigenza universale, una necessità intrinseca alla realizzazione individuale; e quando, dall’altro, la laboriosità diffusa, forgiata dalla rigorosa disciplina imposta dal capitale attraverso generazioni successive, si sarà consolidata come un patrimonio condiviso, un’eredità di una nuova epoca. In breve, quando questa laboriosità universale, nutrita dallo sviluppo delle forze produttive del lavoro, costantemente stimolate dal capitale, animato da un’inesauribile sete di profitto e operante in tali condizioni, raggiungerà la sua piena espressione».Marx scriveva: «In un sistema in cui ciascuno potrà dedicarsi liberamente alle proprie inclinazioni, si verificherà che, da un lato, il mantenimento e l’incremento del patrimonio universale richiederanno alla società solo un numero esiguo di ore di lavoro, e che, dall’altro, la società produttrice adotterà un metodo scientifico per il progressivo sviluppo e la sua riproduzione in quantità sempre maggiori. Di conseguenza, verrà meno quel tipo di lavoro in cui l’uomo si limita a fare ciò che potrebbe delegare alle cose, a vantaggio proprio». Aggiungiamo che ciò segnerebbe la fine dell’intero periodo dell’esistenza e dello sviluppo della divisione del lavoro nella società.

L’antica visione di una fabbrica automatizzata, già delineata nell’era di Marx, e l’auspicio di un’ampia valorizzazione delle risorse naturali attraverso una conoscenza approfondita, a vantaggio dell’intera umanità, rispecchiano le tendenze concrete che convergono verso la creazione delle premesse per l’avvento del comunismo. Questo ideale si concretizzerà progressivamente, man mano che le macchine assumeranno un ruolo sempre più preponderante nello svolgimento dei compiti che possono essere loro affidati, con o senza la supervisione umana, e sempre meno lavoro umano non retribuito risulterà necessario per la perpetuazione della vita sociale. Parallelamente, l’attenzione si sposterà dalla gestione delle persone alla gestione delle risorse materiali. Ogni individuo potrà quindi dedicare una porzione crescente del proprio tempo allo sviluppo libero delle proprie capacità fisiche e intellettuali, ovvero a perseguire ciò che gli procura soddisfazione, a beneficio dell’intera comunità. Questo rappresenta, in definitiva, il fulcro dell’azione comunista.

Il capitale, nella sua funzione storica di relazione sociale imprescindibile per l’evoluzione delle forze produttive, non può pienamente esplicarsi all’interno del sistema capitalistico. Nel suo dominio incontrastato, esso riproduce incessantemente la dimensione materiale del lavoro, impedendo il progresso delle forze produttive ogni qualvolta si manifesti una tendenza a superarla. Di conseguenza, nel capitalismo, la scienza non può concretamente trasformarsi in una forza produttiva autonoma; e, riproducendo il carattere materiale del lavoro, vengono altresì riprodotte le sue forme più arcaiche. L’eliminazione definitiva del capitale in quanto relazione sociale rappresenta la questione cruciale nel periodo di transizione verso la società comunista. In tale fase, è imperativo promuovere attivamente un mutamento nel funzionamento del capitale, attraverso il trasferimento della produzione strettamente materiale al controllo delle macchine, elevando l’attività umana al di sopra della mera dimensione materiale del lavoro.

A tal proposito, è doveroso ricordare che il capitale, parallelamente, alimenta anche una tendenza opposta, che spinge la società verso un ritorno al capitalismo. È contro questa specifica tendenza, sebbene inestricabilmente connessa alla precedente, che la società socialista può e deve opporsi con determinazione, affinché la transizione verso il comunismo possa concretizzarsi. La linea di confine tra il vetusto sistema capitalista e il nuovo ordine comunista si definisce oggettivamente nel punto in cui si presentano tutte le condizioni materiali necessarie per affidare in modo definitivo il lavoro fisico agli strumenti. Laddove tali condizioni non sono ancora pienamente realizzate, a causa di una insufficiente divisione del lavoro e della conseguente mancata creazione delle premesse per l’automazione, il capitale persiste. Nelle aree in cui la completa automazione della produzione materiale di beni è effettivamente possibile, il capitale, in quanto rapporto sociale e forma di lavoro corrispondente, perde la sua ragion d’essere e, di conseguenza, può e deve essere oggettivamente superato.Finché questo non sarà realizzato, il capitale continuerà a riprodursi come una relazione sociale imprescindibile.

La nota questione posta dall’antocomunista medio, “chi si occuperà delle pulizie se tutti faranno ciò che desiderano?”, perde ogni fondamento. La sua retorica si rivela priva di senso non solo perché sussiste una differenza sostanziale tra dedicarsi occasionalmente a un’attività piacevole e poi svolgere un compito sgradevole ma necessario per la comunità, e dover essere costretti a farlo per l’intera esistenza, ma anche perché oggi esistono servizi igienici dotati di sistemi di autopulizia. Ed è solo il predominio della proprietà privata e il fatto che la forza lavoro degli addetti alle pulizie – ovvero, il tempo di vita delle persone – in queste circostanze risulti più vantaggiosa rispetto alle tecnologie, a impedirne una diffusione su larga scala.

In una società comunista, l’individuo, con il suo tempo e il significato della sua esistenza, assurge a valore primario, e ogni cosa acquista rilevanza unicamente in quanto contribuisce a garantire una piena e autentica esperienza di vita per ciascun membro della comunità. L’idea che un essere umano debba impiegare il proprio tempo in attività lavorative ripetitive e prive di significato, attività che potrebbero essere svolte dalle macchine, è inconcepibile in un sistema in cui l’uomo è posto al centro. Questa condizione deve essere superata, e solo allora si potrà parlare di comunismo, fondato su solide basi.

Tuttavia, è doveroso riconoscere che, nella fase di transizione verso il comunismo, per un certo periodo di tempo, persisterà una quota di lavoro non completamente esente da incombenze, finché le tecnologie necessarie per la sua completa eliminazione non saranno sviluppate e ampiamente implementate.L’abolizione del dominio del capitale permetterebbe di ridurre l’orario di lavoro all’attività strettamente necessaria, escludendo quelle mansioni che trovano una logica solo nel contesto di un’economia orientata all’accumulazione, la cui espansione incessante genera risultati che vengono poi sistematicamente distrutti durante le crisi o, più semplicemente, considerati come mere voci di costo. In una fase iniziale, il lavoro necessario – inteso come ancora non completamente liberato – estenderebbe i propri confini. Tale ampliamento risulterebbe imprescindibile per garantire le risorse pubbliche e creare le condizioni per un ulteriore progresso. Tuttavia, l’evoluzione della società verso il comunismo, con il conseguente sviluppo delle forze produttive, implicherebbe una progressiva riduzione di tale attività, fino a ridurne le dimensioni in modo significativo. Oltre una certa soglia, questo tipo di lavoro potrebbe persino essere trascurato, poiché non eserciterebbe più un’influenza determinante sul corso generale dell’esistenza umana.L'elemento cruciale del processo di trasformazione verso il comunismo, di cui si è discusso in precedenza, risiede nel fatto che tale processo, pur non essendo ancora giunto a compimento, mira a creare un tempo libero e a formare un individuo dotato della capacità di agire in modo autonomo e responsabile all'interno della società.
\end{document}
